\section{Elementary full-series bounds}\label{app:scalar}
For $r>0$ set $S_m(r)=\sum_{k=0}^m r^k/k!$.
The positive exponential series and its geometric tail give
\[
 S_m(r)<e^r\le S_m(r)+\frac{r^{m+1}}{(m+1)!}
                  \frac1{1-r/(m+2)}\qquad(0<r<m+2).
\]
The positive logarithm identity is
\begin{equation}\label{eq:atanh}
 \log\frac{1+z}{1-z}=2\sum_{k\ge0}\frac{z^{2k+1}}{2k+1},
 \qquad0<z<1.
\end{equation}
Omitted denominators may be bounded below and powers summed geometrically.
For example
\begin{gather*}
 693/1000<842/1215<\log2<25/36<7/10,\\
 1098/1000<7379/6720<\log3<1+1/12+1/60=11/10,\\
 1609/1000<\log5<13724/8505<81/50.
\end{gather*}
The lower bound for $\log5$ follows from the first eight terms at $z=2/3$,
which exceed respectively
$1.33333,.19753,.05267,.01672,.00578,.00210,.00079,.00030$.
These are exact finite decimals summing to $1.60922$.
The upper bound uses the first three terms and their geometric remainder.
Squaring proves
$\sqrt2>140/99$, $\sqrt3>55/32$, $\sqrt5>324/145$;
these imply \eqref{eq:coefficientcaps}.
Also $\log n/\sqrt n\le2/e<1$, so $a_n\le1$.
The inequality $\Lambda(n)\le n/2$ follows from $\log n<n/2$ for $n\ge2$.

The arctangent addition formula and alternating series give
\[
 \pi>4\sum_{d\in\{2,3\}}\sum_{k=0}^3
       \frac{(-1)^k}{(2k+1)d^{2k+1}}
 =\frac{1538665}{489888}>\frac{157}{50},
\]
where the last difference is $10417/12247200$.
Also
$\frac{22}{7}-\pi=\int_0^1z^4(1-z)^4/(1+z^2)\,dz>0$.
Thus $\pi^2<10$.
Useful Taylor consequences are
\[
 e^{2/5}<3/2,\quad e^{1/5}<5/4,\quad
 e^{8/5}<5,\quad e^3>401/20>20,\quad e^2>73/10.
\]
The bound at $8/5$ follows from the displayed tail formula with $m=6$.

\subsection{The exterior theta envelope}
For $g(t)=(4t^2-6t)e^{-t}$ and $t\ge3$,
\begin{equation}\label{eq:thetatotal}
 \sum_{j\ge2}g(j^2t)<128t^2e^{-4t}.
\end{equation}
Indeed bound each term by $4t^2j^4e^{-tj^2}$.
Consecutive ratios for $j\ge2$ are at most
$(3/2)^4e^{-5t}<1/2$, so twice the first majorant pays the complete tail.
Both $g(t)$ and $t^2e^{-4t}$ decrease for $t>3$.

At $B=4/5$, $S_7(8/5)>99/20$ gives
$t_B>T_0=15543/1000>31/2$.
The logarithm bounds imply $e^B<\sqrt5<9/4$, so $1+e^{-B}>13/9$.
The degree-nine exponential sum at $3$ is
$7280604/362880>401/20$; the degree-four sum at $543/1000$ is
$41297535960801/24000000000000>43/25$.
Hence $e^{15}>20^5(1+5/400)=3240000$ and $e^{T_0}>5572800$.
Since $4T_0^2-6T_0<8731/10$, the leading part of $A(B)$ is below
$78579/724464000$.
The entire tail is below $32768/20^{20}<10^{-12}$, while
\[
 \frac1{9216}-\frac{78579}{724464000}
 =\frac{279936}{9216\cdot724464000}>10^{-8}.
\]
This proves $A<1/9216$.
The logarithmic derivatives of its two summands are at most
\[
 \frac1{1+e^x}+\frac{2(8t-6)}{4t-6}-2t\le6-2t<-25,
 \qquad 4-8t<-25
\]
on $x\ge B$.
Therefore $W(B+h)\le A(B+h)\le Ae^{-25h}$ and
$\int_B^\infty W\le A/25$.

\subsection{A global theta cap and the exterior probability}
Termwise differentiation is justified by locally uniform summable majorants, and gives
\[
 \kappa'(x)=e^{x/2}\sum_{j\ge1}t_jP(t_j)e^{-t_j},
 \qquad P(t)=-8t^2+30t-15.
\]
At zero the first theta term is below $46/49$:
$4\pi^2-6\pi<4(22/7)^2-6(22/7)$, while
$e^\pi>e^{31/10}>22$.
The remaining terms are bounded by $40j^4e^{-3j^2}$,
whose first term is below $1/250$ and successive ratio below $1/10000$.
Thus $\kappa(0)<46/49+40/9999$.
The polynomial $P$ decreases on $[3,\infty)$, with
$P(31/10)=28/25$ and $P(16/5)=-23/25$.
All $j\ge2$ derivative terms are negative.
A positive derivative is therefore possible only on
$0<x<\frac12\log(32/31)<1/62$.
There,
\[
 \kappa'_+(x)<\frac{60}{59}\frac{16}{5}\frac{28}{25}\frac1{22}
 =\frac{2688}{16225}<\frac16.
\]
Hence
\[
 0<\kappa(x)<\frac{46}{49}+\frac{40}{9999}+\frac1{372}
 =\frac{57440653}{60753924}<\frac{19}{20}.
\]
Evenness gives the global assertion and $\sup W<1/2$.
The normalization proved in Appendix~\ref{sec:energy} gives
\[
 \int_0^\infty W
 =\int_0^\infty\frac{\kappa}{2c}\le\frac12\int_0^\infty c\kappa=\frac18.
\]

For completeness, the endpoint cap $\kappa(B)<1/4000$ follows from
$t_B>31/2$, $e^{B/2}<3/2$, $e^{31/2}>5300000$:
\[
 \kappa(B)<\frac{1302}{5300000}
 +\frac{23064}{5300000^4}\frac{10000}{9999}
 <\frac{1302}{5300000}+\frac1{10^9}<\frac1{4000}.
\]
The exponential lower bound is supplied by $S_{24}(31/2)$;
its excess over $5300000$ is
\[
 \frac{37076883949467729399970752217216753}
 {10409396852733332453861621760000}>0.
\]
The tail ratio in this estimate is below $1/10000$.
Furthermore
\[
 -\kappa'-25\kappa
 =e^{x/2}\sum_{j\ge1}t_j(8t_j^2-130t_j+165)e^{-t_j}>0
 \quad(x\ge B),
\]
since at $t=31/2+u$ the polynomial is $8u^2+118u+72$.
It follows that
\[
 \int_B^\infty e^{x/2}\kappa(x)\,dx
 <\frac{(3/2)(1/4000)}{49/2}=\frac3{196000},
 \qquad \mu(O)<\frac3{49000}.
\]

\subsection{All needed interior point estimates}
Since $t(1/10)>19/5>16/5$, the differentiated polynomial proves that
$\kappa$ and $W$ decrease to the right of $1/10$.
Every tail below is the complete bound \eqref{eq:thetatotal}, with the factor $e^{x/2}$ when estimating $\kappa$.

At $x=1/10$,
$t>(157/50)(61/50)>19/5$,
$1+e^{-x}>19/10$, and $e^{19/5}>20(11/5)=44$.
The leading bound for $W$ is $437/1045$; its tail is below
$46208/96000000<1/1000$.
The margin $21/50-437/1045=19/10450$ pays this tail.

At $x=1/4$, $t>(157/50)(79/48)>103/20$,
$e^{1/4}<9/7$ gives $1+e^{-x}>16/9$, and
$e^{103/20}>20(73/10)(29/25)=4234/25$.
The leading bound is $67671/270976$, the full tail is below
$3456/(20^6\cdot7)<1/100000$, and the margin to $1/4$ is
$73/270976$.

At $x=9/20$,
$(157/50)S_3(9/10)=761921/100000>15/2$,
$1+e^{-9/20}>31/20$, and $e^{15/2}>20^2\cdot4=1600$.
The leading bound is $9/124$, the tail is below $7200/20^{10}<1/1000$,
and the margin to $3/40$ is $3/1240$.

For $\kappa(2/5)$, use $t>69/10$, $e^{1/5}<5/4$,
$e^{69/10}>900$.
The leading bound is $207/1000$, and the tail is below
$\frac{19044}{5\cdot900^4}\frac{10000}{9999}<10^{-6}$,
proving $\kappa(2/5)<1/4$.
The first inequality follows from $(157/50)S_3(4/5)>69/10$.

At $x=1/2$ take $t>17/2$, since $(157/50)S_5(1)>17/2$.
Use $e^{1/4}<9/7<4/3$ and $e^{17/2}>20^2\cdot7\cdot(79/48)$.
The leading bound is $952/13825$ and the tail is below
$36992/(3\cdot20^{11})<1/1000$.
The margin $7/100-952/13825=63/55300$ pays the tail.

For $\kappa(3/5)$,
$S_6(6/5)-520/157=88083/12265625>0$ and
$S_5(7/5)-4=1979/46875>0$ give $t>52/5$ and
$e^{52/5}>32000$; also $e^{3/10}<10/7$.
The leading bound is $(10/7)(9256/25)/32000$.
The full tail is below $19840/20^{13}$, by the majorants
$620j^4e^{-10j^2}$ and their ratio below $1/2$.
Their sum is below $1/50$.

At $x=7/10$, take $t>25/2$, using $e^{7/5}>4$.
Also $e^{7/20}<47/33<10/7$ and
$e^{25/2}>20^4(79/48)$.
The leading bound is $33/11060$ and the tail is below
$200000/(7\cdot20^{16})<10^{-6}$.
The margin to $3/1000$ is $9/553000$.

At $x=3/4$,
$(157/50)S_5(3/2)=896627/64000>14$,
$e^{3/8}<3/2$, and $e^{14}>20^4(73/10)=1168000$.
The leading bound is $21/23360$ and the tail is below
$37632/20^{18}<10^{-5}$.
The margin to $1/1000$ is $59/584000$.
All exponential comparisons above follow from the stated positive sums or their geometric remainder formula; no evaluated theta value enters.

\subsection{Gamma constants and all overlap attenuations}
Since $\rho(s)=e^{s/2}/(2\sinh s)$, the inequality
$\rho(s)\le1/(2s)+1/4$ is equivalent to
$2se^{s/2}\le(2+s)\sinh s$.
For odd degree $n$ coefficient comparison reduces to $n\le2^{n-1}$;
for even $n\ge2$ it reduces to $2^{2-n}\le1$.
Thus the power-series difference is nonnegative.
The lower bound $\rho(B)>5/6$ uses $e^{2/5}<3/2$, $e^{8/5}<5$;
the upper bound $\rho(B)<9/10$ uses $e^{2/5}>10/7$, $e^{8/5}>9/2$.

The harmonic error formula gives
\[
 H_n-\log n-\gamma_E
 =\sum_{k=n}^\infty\int_k^{k+1}\left(\frac1x-\frac1{k+1}\right)\,dx
 >\frac1{2(n+1)}.
\]
At $n=5$ and $\log5>8/5$ this yields $\gamma_E<3/5$.
Together with $\log\pi<6/5$, $\pi/2<11/7$ and $\log2<7/10$,
\[
 \cg<\frac65+\frac35+\frac{11}{7}+\frac{21}{10}
 =\frac{383}{70}<\frac{11}{2}.
\]

For the six overlap caps let $\sigma=25/2$.
The reflection exponents exceed $19/8,47/8,35/4$, from the stated logarithm lower bounds.
The positive exponential sums give
\begin{align*}
 e^{19/8}&>7(185/128)>10,\\
 e^{47/8}&>20\cdot7(289/128)>300,\\
 e^{35/4}&>400(73/10)(21/10)=6132>6000.
\end{align*}
For the translations, use
$\sqrt{3/4}<7/8$, $\sqrt{3/5}<4/5$, $\sqrt{4/5}<9/10$, and
\[
 28\cdot3^{12}<4^{12},\qquad
 2000\cdot3^{12}<5^{13},\qquad
 27\cdot4^{12}<2\cdot5^{12}.
\]
These give $(3/4)^\sigma<1/32$, $(3/5)^\sigma<1/500$ and
$(4/5)^\sigma<1/15$, paying every listed pair.
